\begin{statement}{B.29}{Refuted}{Finite-factor lifting does not impose the cusp condition}
Put $B=C_2\amalg C_3$, the free profinite product, with generators
$a,b$ of orders two and three, and cusp subgroup
$I=\overline{\langle ab\rangle}$. There exists a continuous
automorphism of $B$ preserving the conjugacy classes of both
finite factors which does not preserve the conjugacy class of $I$.
Consequently finite-factor automorphism lifting alone does not
settle the peripheral profinite lifting step in B.28.

The following corrected lifting statement is true. For every
continuous epimorphism $q:B\to Q$ of profinite groups, every
continuous automorphism $\alpha$ of $Q$ preserving the conjugacy
classes of $q(C_2)$ and $q(C_3)$ lifts to a continuous automorphism
$\beta$ of $B$ preserving the conjugacy classes of $C_2,C_3$,
with $q\beta=\alpha q$. No assertion about $I$ is included.
\end{statement}
\begin{proof}
We derive the corrected assertion from
\citet[Main Theorem, pp.~2217--2219]{B_Efrat1997}.
That theorem says that if a finite group $G$ is generated by
subgroups $G_1,G_2$, and its quotient $H$ is generated by prescribed
conjugates of their images, then $G$ is generated by conjugates
of $G_1,G_2$ with those prescribed images. We give the full
passage to the automorphism assertion, because it is exactly
where ordinary generator lifting can lose a condition.

For any finite quotient $Q_0$ of $Q$, and any finite quotient
$B_0$ of $B$, take the joint finite quotient $G$ of $B$ in
$B_0\times Q_0$. Its images of the two factors generate it.
Apply Efrat's theorem to $G\twoheadrightarrow Q_0$, prescribing
the subgroups generated by the images of $\alpha q(a)$ and
$\alpha q(b)$. The resulting conjugates generate $G$.
An element of order two has one generator; an element of order
three has two. Thus, after choosing $e\in\{1,-1\}$, the pair
has the form
\[
          uau^{-1},\qquad vb^ev^{-1}
\]
in $G$, and maps to the specified individual elements in $Q_0$.
Choose lifts of $u,v$ to $B$. The same assertion holds if a
factor has trivial image in $Q_0$.

In the compact space $B\times B\times\{1,-1\}$ impose the
closed conditions that the two displayed elements have their
specified images in $Q_0$ and generate $B_0$. Any finite list
of these conditions is included in one joint finite quotient
of the same kind, so the conditions have the finite intersection
property. Compactness gives a pair satisfying them for every
$Q_0$ and $B_0$. The universal property of the free profinite
product supplies an endomorphism $\beta$ of $B$. Its image
surjects onto every finite quotient of $B$, hence is dense;
it is also compact, hence closed, so $\beta$ is surjective.

A topologically finitely generated profinite group is Hopfian:
for each fixed index it has finitely many open normal subgroups,
and inverse image under a surjective endomorphism is an injection
on this finite set, hence a permutation. Its kernel consequently
lies in every open normal subgroup and is trivial. Applying
this argument to $B$ proves that $\beta$ is an automorphism.
The conditions on every $Q_0$ give $q\beta=\alpha q$.
This proves the corrected lifting assertion with its topology
and comparison map specified.

For the counterexample we use the same finite-subgroup theorem
with prescribed generating elements of one finite quotient;
they need not be obtained from an automorphism of that quotient.
In $G=\operatorname{PSL}_2(\mathbb F_{13})$ put
\[
 A=\begin{pmatrix}0&1\\-1&0\end{pmatrix},\quad
 B_0=\begin{pmatrix}0&1\\-1&-1\end{pmatrix},\quad
 H=\begin{pmatrix}0&1\\-1&2\end{pmatrix}.
\]
All matrices are taken modulo scalar sign. Direct multiplication
gives
\[
 B_1=H^{-1}B_0H=
      \begin{pmatrix}3&6\\1&11\end{pmatrix},
 \qquad
 |A|=2,\quad |B_0|=|B_1|=3,
 \qquad |AB_0|=13,\quad |AB_1|=7.
\]
Both $(A,B_0)$ and $(A,B_1)$ generate all of $G$.
The executable \texttt{b\_factor\_lifting.py} checks these
claims by exact matrix multiplication, enumerates the full
determinant-one quotient of order $1092$, and compares both
generated subgroups with it. Its JSON and printed output
contain the complete small matrix certificate.

Let $q:B\twoheadrightarrow G$ send $(a,b)$ to $(A,B_0)$.
For each joint finite quotient of $B$ refining $G$, Efrat's
theorem supplies generating conjugates mapping to the subgroups
$\langle A\rangle,\langle B_1\rangle$. Choose the order-three
generator as above. The identical compactness and Hopfian
argument gives an automorphism $\beta$ with
\[
                    q\beta(a)=A,\qquad q\beta(b)=B_1.
\]
If $\beta(I)$ were conjugate to $I$, the image of its topological
generator $\beta(ab)$ in $G$ would have the same order as
$q(ab)$, namely $13$. Its actual order is $7$. This is a
contradiction. Thus the cusp condition is independent of this
finite-factor lifting principle.

Applied to B.28, the principle can lift the action on the two
elliptic factors of the relative modular group, but its conclusion
does not retain the third, cuspidal condition or the normalized
\S47 word equation. The explicit counterexample identifies the
missing condition; it is not presented as a nonpentagonal
peripheral profinite lift.
\end{proof}
